

\section{Free Connection among Free Congregations}

Opening Address at the Free Church’s First Meeting,
by Professor Sven Oftedal.

Report of the Annual Meeting of the Lutheran Free Church, 1897, p. 62ff.

It may well be maintained as one of the so-called self-evident truths, a matter taken for granted, that if the congregations are free, then the mutual connection among them must also become a free connection.

Nevertheless, or rather precisely because it is self-evident, it is so exceedingly difficult, not to say impossible, to carry out in practice.

It is self-evident that a sinner, because he is a criminal, can be saved only by being pardoned; and yet how difficult, indeed impossible! And how many other ways has not human reason and theology devised in order to avoid this one narrow way, to be saved by grace alone!

Free connection among free congregations is extraordinarily difficult to realize:

* 1. On account of the prerequisite.
* 2. On account of the apparent self-contradiction in its execution.

The prerequisite is, namely, free congregations, or a true freedom of the congregation.

By free congregations we understand, in this connection, not congregations that are freed from the supremacy of the state, or even from the tutelage of the priesthood or of the papacy, but congregations that are in truth free in Christ, set free in their individual members by the Gospel, pervaded, governed, and driven by the Spirit of God according to the word: "Where the Spirit of the Lord is, there is freedom" (2 Cor. 3:17).

If congregations of that sort are the prerequisite for the aim here under consideration, then anyone who knows anything of the condition of our congregations will at once understand not only how great, but how difficult, the task lying before the Free Church is.

Only in the same measure as the congregations are awakened, made alive, and set free by the Holy Spirit and power of God can there gradually be any question of establishing a free connection among them. The Free Church therefore is only as it continually becomes.

But even if the prerequisite is present, the execution is equally difficult, because it involves an apparent contradiction. God’s congregation is, namely, at the same time that it is many, yet always only one. We believe one holy catholic Church. Connection among the congregations always aims to present an outward visible form of this unity.

But here we are met by the fact that everything in the world is "in part": "we understand in part and we prophesy in part; and only when the perfect comes shall that which is in part be done away" (1 Cor. 13:9-10). Therefore freedom also, true spiritual freedom, is always in part; and naturally every free connection is also in part; it too is only as it becomes, and therefore does not satisfy the demanding reason’s claim for an outward, sensible, visible, and perfect unity. The Church cannot at one and the same time be both many and one, says reason.

"Free connection among free congregations," as a presentation of the unity of the Church, is therefore only a beautiful thought-experiment, an unattainable ideal, people say; it is not only a difficulty, but an outright impossibility.

The history of the Church therefore shows us that the "Fathers" have always had to seek more practical and rational ways by which the connection among the congregations might realize the great, inescapable goal: the unity of the Church.

These rational methods of unity, no doubt often forced forth by historical conditions, may be gathered under three great groups:

* 1. The episcopal principle.
* 2. The territorial principle.
* 3. The synodal principle.

By all these methods one succeeds in presenting an outward, unmistakable form of the unity of the Church, but it always takes place at the expense of the congregation’s freedom and individuality. They all run aground on the contradiction: many and yet one; one and yet many (cf. 1 Cor. 10:17).

The episcopal principle, which in reality makes the congregation identical with the bishop, leads in its consistent development to the papacy; that is to say, to the unity of the Church represented by the highest bishop, Christ’s vicar. Although the Catholic Church, with its personal visible head and its fully developed organization, has incontestably gone farthest in presenting the unity of the Church on an admirable and almost world-embracing scale, yet we Lutherans acknowledge that this has taken place at the expense both of doctrine and of the congregation. The doctrine of God’s Word has been corrupted, and the congregation almost altogether obliterated. The priest has become everything, the congregation nothing. The existence of our Lutheran Church is precisely a serious and definite protest, grounded in the Word of God, against this principle.

The state-church principle, which is also called the territorial principle because the boundaries of the congregation are determined by geographical divisions or territories, has early been tried both in Catholic and in Protestant lands. Every time it has been tried, it has led to complete worldliness. We have seen its fruits in our old fatherland, where a hundred years ago almost every single priest was a rationalist or freethinker, and Hauge, on account of his revival work, was cast into prison. When God and the world are apparently reconciled, it is God who is excluded; when the congregation and the state are reconciled or united, then there does indeed come about as complete a union as human power and compulsion can produce, but also in this form of the unity of God’s Church it is the worldly state that becomes everything; the congregation becomes nothing but an outward form; it is in reality dissolved according to its spiritual essence and passes over into becoming equivalent to a district or a municipality, a township or county, in short a purely worldly organization with the Word and the Sacraments as a label pasted on from the outside.

The synodal principle has also, from the earliest time and right down to our own days, been tried in both Catholic and Protestant lands. It is what in this country we are most familiar with under the name of the associational principle. It has great advantages over the two preceding. In theory, and even in constitutions, it can strongly emphasize and acknowledge the independence of the individual congregation. But since this principle too aims at an outward, visible, and organized unity of the Church and seeks to realize this in the "synod" or the "association," a consistent development seems unavoidably to lead, little by little, to setting up the "synod" or the "association" as a higher sort of congregation, a form of unity of all the congregations belonging to the association. Thereby the congregation has already lost its place as the only outward form appointed by God for His kingdom on earth; it has been reduced to being an intermediate stage leading to a higher form of God’s Church, fashioned by men.

Once the development has entered upon this slippery slope, though in many gradations, congregations will unceasingly have to withdraw more and more, until at last they are even reduced to a mere supporting foundation for the financial operations of the "association." The synod, as the totality and unity of all the congregations belonging under it, will in the course of development come to present the outward form of this higher unity in the priests, who will then gradually and involuntarily regard themselves individually as special members of the congregation, placed higher by ordination, and in their totality, as an "estate" or "priesthood," consider themselves the actual representatives of the Church and of its unity. It is here that the synodal principle in its development will meet the papacy, no doubt on a smaller scale, but on no less effective a one. Here too the congregation is sacrificed for outward priestly church-unity.

These different ways of presenting the unity of the congregation, in order to avoid the self-evident and simple, but for human reason impossible thing: the preservation of freedom both for the individual congregation and for their mutual connection, have now for nearly nineteen hundred years again and again been tried and weighed and each time found wanting, because in the course of their development they have all blurred or corrupted "the pillar and ground of the truth, the congregation of the living God," and have set up their own organizations in its place, partly priestly, partly civil.

For serious-minded, independent Christians, therefore, and especially for us, who have to build a Lutheran congregation as if anew in this country, there remains nothing else than to attempt the "impossible," according to the word that stands written: What is impossible with men is possible with God. "God’s temple, His congregation, is holy," says Scripture, and adds: "which you are." The congregation is not something human; it is God’s work; it can therefore neither be built nor united in its multiplicity into one, whether by might or by power, but only by the Spirit of God. He alone is the one who can establish a "free connection among free congregations."

It is true that the history of the Church after the apostolic period shows us that the Church’s foremost and most earnest leaders, the Fathers, down through the ages, have almost all regarded one of the three above-named principles as divinely warranted in the work of presenting the unity of the Church outwardly; and with apparent justice the old question may also be raised here: "Are you wiser than all our Fathers?"

But first, the same question was put to Luther. Therefore it is true Lutheran doctrine that one Word of God weighs more than all the Fathers. Both the congregation’s doctrine and its order, two things which at many points are inseparable, must be tested by the Word of God and by the Word of God alone. Fathers, history, tradition may be of inestimable value for illumination and guidance in spiritual matters. But when it concerns the very pillar and ground of the truth, the congregation of the living God, the body of Christ, its essence and unity, then it is our duty and our right, most of all as Lutheran Christians, to pass beyond both Fathers and Tradition and to seek in the Word itself the knowledge and strength that we need.

Second, the New Testament gives all desirable instruction about the congregation both according to its divine essence and according to its historical reality.

There now therefore present themselves two questions:

* 1. Was the congregation, as it meets us within the period of the New Testament, of the same kind and essence as the congregation now?
* 2. Was there among these New Testament free congregations any unity of connection, and can a similar one be brought about among our congregations?





















% Page 64





\subsection{The New Testament Congregation}

The first fact that meets us in this connection is that the New Testament speaks only of “the congregation,” the individual congregation; in catholic and state-church usage theological language has, in place of “the congregation,” put “the Church,” a word which by its very concept has most often — and so also in the free-church clerical bodies here in question — swallowed up the congregation and usurped its place and right. From this standpoint the congregation is therefore regarded essentially only as an intermediate link toward the higher unity, “the Church,” and is accordingly, in distinction from “the Church,” commonly called the district congregation or the local congregation.

We find, further, that before the congregation was created on the day of Pentecost and set into the world, in other words in the Gospels, this word, “the congregation,” occurs only twice, and that in Matthew. First in Matt. 16:18, where the Lord says to Simon: “You are Peter, and upon this rock I will build My congregation,” the passage upon which the whole Catholic church theory has been built, but against whose Catholic interpretation all Protestants have objected with fully valid reason. Next in Matt. 18:17, concerning discipline in the congregation: “— — tell it to the congregation, and if he does not hear the congregation,” etc. If the Lord here does not refer simply to the Jewish synagogue, there is at all events no doubt that what is meant here is the congregation in the same sense in which the word is otherwise used in the New Testament, the individual local congregation, large or small.

All the more frequently is the word congregation used in the writings after Pentecost, in all more than one hundred times. More than half of these fall in four writings, namely the Acts of the Apostles, which may rightly be called the history of the New Testament Church, the two Epistles to the Corinthians, and the Epistle to the Ephesians, the very writings which also give us the clearest insight into the nature of the congregation according to the presentation of God’s Word.

Among all these many passages where “the congregation” or “the congregations” is meant, there is not one that can be explained of “the Church” in the sense of an outward church-body consisting of some or all congregations as a higher outward unity of them. God’s Word does indeed present all congregations as a unity, as one body — (Eph. 1:22 and 4:16; Col. 3:11; 1 Cor. 10:17, etc.) namely Christ; but it is a spiritual unity, according to that word, that Christ is all in all, that those who are baptized into Christ have been clothed with Christ, and that there is neither slave nor free, male nor female, Greek nor Jew. And what is said of this unity of all congregations is also said, and in the same passages, of the individual outward congregation, namely that it is Christ’s body, or that it is the body and Christ the Head, and that every individual member in this outward congregation is a member on Christ’s body.

In God’s Word there is lacking every indication that for the outward connection between these congregations there was any arrangement or any common government which made necessary an outward head over the congregations in their totality as a higher form of the congregation. Not even the apostles had, so far as appears from God’s Word, any authority in other congregations than the first congregation in Jerusalem and possibly the first congregations in Samaria. Indeed, even concerning the apostles’ activity there is in God’s Word a remarkable — and yet in God’s merciful wisdom well-grounded — silence, in order that there might be in the congregations the least possible temptation to set men over Christ. The New Testament even speaks disproportionately much more of the thirteenth apostle, “one untimely born,” “the least of the apostles,” “who was not worthy to be called an apostle” (1 Cor. 15:9), than of “the very chiefest apostles” (2 Cor. 12:11), although he, as he says, “is nothing,” evidently in order to show the congregations that God’s grace is not bound to human persons, but to Christ alone, and to counteract the tendency which already in the first congregations, and especially in the spiritual congregation at Corinth, showed itself so strongly, namely party-spirit: that one held to Cephas, one to Paul, one to Apollos, one to Christ, and “you endure it if anyone makes slaves of you, if anyone devours you, if anyone takes you captive, if anyone exalts himself, if anyone strikes you in the face!” (2 Cor. 11:20).

It is undeniable that it is in essential part out of this fleshly craving in every human being to regard persons and to find something outward to build upon, that there developed quite early in church history a strong urge to lift the apostles and the first Christian workers so high up that Christ and His body, the congregation, became in the same proportion of subordinate significance; and this not only in order to build monuments over the prophets, but chiefly thereby to secure for themselves, as the leaders and governors of the “congregation” and later of the “Church,” as priests and bishops, a recognition and authority which they thought they could uphold only as “the successors of the apostles,” “Peter’s heirs,” “Christ’s vicars,” etc., the principle upon which the whole Catholic church order, all state churches, and all clerical bodies are in reality built.

A wholly different picture presents itself when we consider the congregation, without regard to the subsequent history and in part very fleshly development, as God’s Word presents it to us, whether the congregations are considered in and for themselves or in their mutual relation to one another.

There are here portions of God’s Word which are very little heeded because they do not fit into dogmatics and into presuppositions and traditions belonging to a much later time.

\subsection{The First Congregation of God}

in Jerusalem was, as we all know, created or born by a miracle of the same kind as the creation of the world and Christ’s becoming man, in that God’s Spirit immediately filled the twelve apostles and the others, in all about one hundred and twenty persons, men and women, who were assembled on the day of Pentecost and awaited “the promise of the Father,” “the baptism with the Holy Spirit” (Acts 1:4–5), and received the power of the Holy Spirit which came upon them. Thus the congregation, according to its first origin, is divine, and by the Spirit and the Word has become the firstfruits of God’s saving order in the world, a leaven which shall work through to the ends of the world, everywhere alike holy by reason of the means of grace, according to that word, that if the first of the dough is holy, then the whole dough is likewise, and if the root is holy, then the branches are likewise (Rom. 11:16).

Concerning the growth and development of this first congregation it therefore stands written that when Peter had spoken, there was a multitude which received the Word and was baptized, and about three thousand souls were added to the congregation that same day. Thus the first congregation advanced both in number and in the power of God’s Spirit, and the multitude of the believers were all of one heart and one soul, and all were filled with the Holy Spirit and spoke the Word of God with boldness (Acts 4:31–32).

And yet even then the congregation was not mere spirit. It consisted of human beings with flesh and blood, who formed an outward association or fellowship, had to meet in an earthly meeting-house, had to do with God’s things and worldly things, indeed even with food. What is unavoidable in every regenerate person, the conflict between flesh and Spirit, therefore soon made itself felt also in this glorious congregation filled with the Spirit, and brought it not only into the sufferings that come from the world’s hatred and resistance, but also into those inward conflicts whereby there is filled up what is lacking in Christ’s afflictions for His body (Col. 1:24).

There were Ananias and Sapphira — the leaven of the Pharisees, which is hypocrisy — and there were jealousy and strife between the Greek-born and the Hebrews over — the food.

No congregation any more than any individual Christian can, so long as it is an outward visible fellowship upon earth, avoid these inward conflicts between the flesh and the Spirit; but if this necessity is not concealed or covered over by fleshly sloth or self-indulgence, then God will also by His Spirit unceasingly give the congregation those powers and abilities whereby it can gain the victory in the conflict and learn by experience that the afflictions work in a refining and purifying way, proven character, patience, and hope (Rom. 5:3).

So it also went here; the congregation in Jerusalem, by the apostles’ counsel, chose out of its own midst men equipped by God to meet and overcome the threatening danger, and after the first deacons were chosen it stands written that God’s Word advanced, that the congregation multiplied, and that even many priests were converted.

Founding of Congregations

Neither was it the apostles who founded new congregations. So far as one can see from God’s Word, the apostles took no measures to carry the congregation farther beyond Jerusalem; they had more than enough to do there.

It was through a particular and unexpected leading of God that new congregations were founded. After Stephen’s murder a persecution broke out, after which the whole congregation in Jerusalem was driven out over Judea and Samaria except — the apostles. Of these scattered members of the congregation it stands written that they “proclaimed the word of the Gospel wherever they came.” Thus, through the deacon Philip, the congregation in Samaria was founded. This, that “Samaria had received the Word of God,” awakened to such a degree the attention of the apostles who remained in Jerusalem, and no doubt also their astonishment, that they sent Peter and John down to see how the matter stood, and they prayed that the converts might receive the Holy Spirit through the laying on of hands, which they also received, so that in that state they all became priests under the new covenant (Acts 8:14 ff.).

In the same way Philip became a means by which the Gospel and the congregation came to Ethiopia. And not that alone, but the scattered believers brought the Word and thereby the congregation to the surrounding lands, Phoenicia, Cyrene, and Syria (Acts 11:19–20), and many other places where we later find congregations without God’s Word specifically mentioning by whom it was brought there, as for example Alexandria and Rome.

But there was here still another hindrance, in part among the apostles and the congregation in Jerusalem itself; namely, that the separation between Jews and Gentiles was regarded as so great that in their opinion there could be no question of preaching the Gospel directly to the Gentiles, or of there being any hope that Gentiles could receive the Holy Spirit and thus truly become a congregation of Christ.

Here too the Lord had to use, not the apostles’ exaggerated wisdom or authority, but quite other and in part extraordinary means in order to make both the apostles and the Jerusalem congregation understand that Jesus died and rose again for all, and that He therefore both can and will save all — Jew or Greek — and make them members of His body, which is the congregation.

First, there were men from Cyprus and Cyrene who first spoke to Greeks in Antioch and there formed a great congregation of believing Gentiles. At this again the congregation in Jerusalem was astonished and sent Barnabas to investigate the matter, which he found altogether in good order, since he, as it is written, was a good man (Acts 11:20 ff.).

Second, the Lord compelled Peter himself, despite all his resistance, to go to the Gentile Cornelius, a step for which he had to justify himself before the congregation in Jerusalem (Acts 11:1 ff.).

Finally, however, God had appointed a particular instrument for the founding of congregations among the Gentiles, the last man any human being could have thought of, a persecutor, Paul, who therefore also received the grace to labour more than all the others, and received so deep an insight into this great mystery, the congregation, Christ in us, as no other did (Eph. 3:4).

And after these preparations the congregation then went on its victorious course among the nations, in that indeed Paul and his fellow workers laboured chiefly among the Gentiles, Peter and the other apostles chiefly among the Jews, but all who themselves had received a share in Christ through faith and baptism zealously proclaimed everywhere, wherever they went, that salvation from sin and death which is found in the Gospel of Jesus Christ, as we see in such as Aquila and Priscilla, Apollos, and Epaphroditus. Where God’s Word fell down into any heart, there it became a seed which grew up and bore fruit, and the fruit was congregations here and there, everywhere in the Roman Empire, as mentioned above, and this without there being any common government or bond between them except God’s Spirit and God’s Word, “the one Spirit, the one calling, the one Lord, faith and baptism, the one Father,” as it is written, “who is over you all and through you all and in you all” (Eph. 4:4 ff.).




% Page 71




The Composition of the Congregations

and their further inner development in the apostolic age also emerges clearly enough from the New Testament writings. At first they consisted of Jews alone, afterward also, and in many places exclusively, of Gentiles from every nation. But whether the congregations were founded among Jews or Gentiles, they consisted of believing men and women and nothing else. Indeed, from the beginning there was no reason whatever why anyone should join the congregation unless they had been delivered from the burden of sin and the fruit of hell through faith in Jesus’ blood and therefore were ready to bear His cross. Therefore, too, in the beginning they were called “brothers,” “disciples,” and “saints,” and only received the name Christians as a term of reproach. But all together they were called, from the first moment until now, the congregation.

But as Christian experience teaches us to this very day, both with respect to individuals and to congregations, not all that is sown bears fruit with perseverance. What is begun is not ended, and Satan is never more active than where the life of Christ manifests itself powerfully. The leaven of the Pharisees and the leaven of the Sadducees soon began their baleful work in the first congregations. And thus, as we actually find the congregations, even within the period of the New Testament itself, they were without exception all infected by the spirit of carnality and very far from consisting exclusively of simple and living Christians.

To hide this fact from oneself proceeds either from complete ignorance of God’s Word in this matter, or from hypocrisy and falsehood, which through centuries of repetition have almost attained prescriptive right as a legitimate doctrine. People have grown accustomed to looking upon the first congregations, with their fullness of the Spirit and boldness of faith, as altogether perfect, an unattainable ideal, a company made up of nothing but saints and martyrs, a peculiar divine wonder-time, and so forth; and then, with apparent right, they have concluded: “We cannot expect to be like the first Christians.” Thereby they have found an excuse for the crying contrast between the manifold dead mass-congregations among us and the apostolic congregations, and in part have further quieted themselves with the fleshly, self-satisfied thought that “it was another matter in former days,” and “that we must be thankful that our congregations are as good as they are”; whereas the truth is this, that if the congregation is the work of God’s Spirit, then our congregations too ought at least to become just as living and vigorous, just as bold in faith and patient, just as self-sacrificing and missionary as in the first Christian age.

For what do we find in this first apostolic age? First, Ananias and Sapphira and quarrels about food, as has already been mentioned. Next, that false teachings, grounded in the current Phariseeism and the fleshly unwillingness to let oneself be saved by grace alone, to be justified by faith without the works of the Law, quite early stole into the congregation in Jerusalem and outside it.

In the Asiatic congregations this tendency was carried much further, in that they set up a peculiar worship of God (the worship of angels) and an exaggerated asceticism or self-denial, which in their opinion was especially meritorious, in place of faith and grace alone (Col. 2:18ff.; 1 Tim. 4:1ff.). On the other side, in many places people began to set worldly wisdom, long genealogies, and profane fables in place of the simple Gospel (1 Tim. 4:7; 6:20), to deny or explain away the resurrection (2 Tim. 2:18; 1 Cor. 15:12), to explain the Lord’s return for fleshly advantage (2 Thess. 3:6), and to use Christian freedom for gain or as an occasion for the flesh (1 Tim. 6:5).

But not only that. Already early the spirit of falsehood and the lust to rule began to creep into the congregations. In the first Thessalonian epistle, which is the earliest letter we have from Paul, he warns against such as either spoke in his name or even sent false letters in his name (2 Thess. 2:2), for which reason at the end of the epistle he expressly calls attention to the way in which he writes his name. In Asia, he says that “all had turned away from him,” and besides this he mentions explicitly two persons, Phygellus and Hermogenes. Demas forsook him out of love for the world, and Alexander the coppersmith had done him much evil (2 Tim. 4:10 and 14). Compare with this 2 Tim. 3:1ff.; likewise 2 Peter and the Epistle of Jude. John says that Diotrephes, who likes to be among the foremost, would not even acknowledge him, but refused to receive the brothers and cast them out of the congregation! (3 John 9ff.) And just as Jesus had done with His disciples, so Peter finds it necessary (1 Pet. 5:1ff.) to admonish the elders both against serving for the sake of gain and against lording it over the flock.

And this is not all. In almost every single epistle there are warnings not only against the above-mentioned perhaps somewhat finer spiritual sins, but also against all manner of gross vices, such as were readily carried over from a surrounding heathen world corrupted to the utmost, to which the members of the congregation had every temptation to fall back. (Eph. 5:3ff.; Col. 3:5ff.; 1 Thess. 4:3ff.; Gal. 5:13ff.; etc.)

For to Paul and the other New Testament writers the congregation does not stand forth as something loose and indefinite, something invisible and general, believers wherever they might be found, without regard to their outward visible union and assembly; rather, it stands forth as a distinct circle consisting of men and women, slaves and free, rich and poor, families and children, separated from the rest of the world by this their outward ordering and by the Word and the Spirit by which they had been separated, and in which nevertheless Satan had been able to produce all the deviations and defections and sins mentioned above.

In almost all the epistles Paul sends greetings to named persons in the congregation, and sends greetings not only from individual persons but from whole congregations, just as he also mentions Gospel work as proceeding from the individual congregations as a whole. In the Epistle to the congregation in Rome, where he himself had not personally been, he nevertheless mentions in his greeting by name so many that their number, including the greeted families, can scarcely be put at less than fifty persons. And if he had such knowledge of the individual members of a relatively unfamiliar congregation, how closely must he not have known the individual members in the congregations he himself had founded, such as, for example,

The Corinthian Congregation?

We will therefore dwell briefly on this.
We have preserved for us in God’s Word two epistles from Paul to this congregation, for which he expresses such heartfelt love and such deep concern.

What, then, was the condition in this congregation?

There were many and great faults. First factionalism: one held to Paul, one to Cephas, and so forth. There were gross sins of immorality, and worse still, a lack of that spirit of holiness which exercises discipline against open ungodliness. There were greed and litigiousness, so common both among Jews and Gentiles. There was carelessness with respect to the holy Lord’s Supper, and class-pride, so that one held himself better than another, drank himself drunk at the meals, and let the poor brothers go hungry. There was spiritual pride over their great spiritual gifts, for which reason they also esteemed the more conspicuous gifts, such as tongues, above the simpler but more necessary gifts, prophecy, mutual brotherly admonition, and above all love. False fleshly teachings about the resurrection had crept in, and lukewarmness in Christian works of mercy. They allowed themselves to be led astray by false intruding brothers, who carried on intrigues against Paul and sought to destroy the fruit of his labor. Therefore he speaks sharp words of punishment over them and says: “I fear that, when I come, perhaps I shall not find you such as I wish, and that I shall be found by you such as you do not wish: that there be strife, jealousy, wrath, quarrels, slanders, whisperings, puffed-up pride, disorders; that when I come again, my God shall humble me among you, and I shall mourn over many who formerly sinned and have not repented of the uncleanness and fornication and sensuality which they have committed” (2 Cor. 12:20-21).

“I fear that, as the serpent deceived Eve by his cunning, so your minds shall be corrupted and led away from faithfulness to Christ” (2 Cor. 11:3). “Some of you have no knowledge of God; I say this to your shame” (1 Cor. 15:34), “and many among you are weak and sickly, and a whole multitude sleep” (1 Cor. 11:30).

And, as in Corinth, so there seems more or less, to judge from the exhortations of the apostles and from the seven epistles in Revelation, to have been great deficiencies in all the other congregations as well.

And yet Paul writes to this same Corinthian congregation: “I thank my God always for you, for the grace of God which is given you in Christ Jesus, that in Him you have been made rich in all things, in all utterance and knowledge, even as the testimony of Christ has been confirmed in you, so that you come behind in no gift of grace, as you wait for the revelation of our Lord Jesus Christ” (1 Cor. 1:4ff.).

And in the second Epistle to the Corinthians he says: “Make room for us! We have wronged no one, corrupted no one, taken advantage of no one. I do not say this to condemn you; for I said before that you are in our hearts, to die together and to live together. I have great confidence in you, I boast much of you, I am filled with comfort, I have an exceedingly great joy in all our affliction. For though I grieved you with that epistle, I do not regret it, though I did regret it; for I see that the epistle, though only for a time, grieved you. Now I rejoice, not that you were grieved, but that you were grieved unto repentance; for you were grieved according to God, so that you might in no way suffer loss through us. Therefore we have been comforted by your comfort; but we rejoiced far more over Titus’ joy, because his spirit has been refreshed by you all. For if I have boasted much of you to him, I have not been put to shame; but just as we spoke all things in truth to you, so also our boasting before Titus has become truth, and he has an exceedingly great love for you when he remembers the obedience of you all, how you received him with fear and trembling.”

Such was the composition of the first apostolic congregations. There were vessels unto honor and vessels unto dishonor. They were like a Christian still yet in the world, surrounded and permeated by its sin and entanglement; they were neither perfect, nor had they attained it. But they pressed on after it; they had the life of God’s Spirit in them; they had within themselves an unceasing witness concerning sin and righteousness and judgment, so that they were not allowed to sink either into security in sin, or into comfort drawn from pure doctrine or from holy living, but kept themselves awake unto a daily sorrow and daily repentance unto God.

And thus they were enabled to carry on the work of the Gospel in the world, though in weakness, yet in inward sincerity, in the power and love of Christ’s Spirit; and although the congregations formed neither an episcopal church nor a priestly order, yet there was among the congregations in their mutual relation a manifestation or realization of the unity of the Church such as has perhaps scarcely ever existed since.


% Page 78



\subsection{The Congregations’ Free Connection}

Just as the congregations in the apostolic age were, in the full sense of the word, free congregations, so also was their mutual connection. God’s Spirit and the love of Christ were that which held them together in unity around the common labor of bringing the Gospel to the ends of the earth.

Therefore, in the New Testament the word congregation is used both to denote the individual local congregation and the unity of all the congregations, that which in our language is ordinarily expressed by the word “the Church,” the one holy catholic Church.

This is already beforehand an expression of the fact that all the congregations, according to their inward spiritual being, by the common doctrine of salvation and the common life of love, felt themselves to be one, and therefore would also find ways and means, as circumstances required or allowed, to give this unity a practical form in their relation one to another.

The first impulse was given by the persecution after Stephen’s murder.

Those who were scattered abroad (except the Apostles) preached the Gospel everywhere they came. That is the right spirit. The fruit was a considerable number of congregations, round about in Samaria, Galilee, Phoenicia, Syria, Cilicia, Cyprus, Cyrene, and no doubt also Egypt, Italy, and Babylon (see Acts 9:31 and Acts 11:10).

Of some of these congregations we hear later in passing, for example in Tyre and Caesarea and Ptolemais (Acts 21:4), but otherwise the Word of God contains no register of founded congregations; there was as yet no “church statistics.”

Yet their cooperation and inward connection were not for that reason any the less lively; both were promoted and guided by the persecutions. Of the means of connection the following may be named:

1. Messengers or Emissaries and Visitors

The most remarkable example of this is the congregation in Antioch, the first Gentile congregation. It sent out Paul and Barnabas on their first missionary journey together with Mark. The former therefore seems to have regarded this congregation as his headquarters, although it is not expressly mentioned that he was sent out more than that one time. As the Apostle to the Gentiles in a special sense, he afterwards journeyed about as the Spirit drove him, visited the congregations already founded and established new ones, and, as it seems, everywhere took fellow laborers and other messengers with him from the founded congregations (Silas, Timothy, Titus, Luke, Demas [Col. 4:14], Sosthenes, Gaius, and Aristarchus [Acts 19:29], etc.).

But what took place in Antioch we find clear traces of as having taken place in other congregations, although no special description of it is given in each particular case.

Thus Barnabas, after the dispute with Paul, continued the work together with Mark, it being mentioned that he went to Cyprus, without its being further stated where the journey afterwards led. The same Mark we find again considerably later in Babylon together with Peter (1 Pet. 5:13), and elsewhere with Timothy (2 Tim. 4:9).

From Jerusalem many envoys went out, not always of the best sort, round about through the lands, also where Paul had founded congregations, and in part even worked against him (Acts 15:24).

Aquila and Priscilla we find both in Corinth, Ephesus, and Rome, everywhere laboring for their Savior, so that they not only assist Paul, and win thanks from all the churches of the Gentiles (1 Cor. 16:19), but also become the means of bringing Apollos all the way to a living faith in Christ and of making him a traveling instrument for the Gospel.

From Cenchreae the deaconess Phoebe travels to Rome and is commended by Paul.

In his greeting to the Romans Paul mentions five fellow laborers, among them a woman, Mary.

Paul sends Titus and Timothy and Tychicus and others from congregation to congregation.

In 1 Cor. 16:15 he mentions the household of Stephanas as one that has especially devoted itself to serving the saints, wherefore hospitality also is warmly recommended in so many places (Heb. 13:2; Rom. 12:13; 1 Tim. 4:9).

The congregation in Berea sent brethren to accompany Paul to Athens (Acts 17:14).

The Jerusalem congregation sent Judas and Barnabas with a letter to Antioch (Acts 13:21).

The congregations sent Paul with aid to Jerusalem (Acts 20:4; cf. 2 Cor., chs. 8 and 9).

Agabus came down from Judea (Acts 21:10); and in 11:28 Antioch sends Paul and Barnabas with aid to Jerusalem.

Disciples from Caesarea traveled with Paul from there in order to procure him lodging in Jerusalem with Mnason, an old disciple from Cyprus (Acts 21:15).

Disciples in Puteoli received Paul, and brethren from Rome met him (Acts 28:14), just as the great companies followed him with wives and children from Tyre (Acts 21:5).

The congregation in Corinth sent Stephanas, Fortunatus, and Achaicus to Paul in Ephesus (1 Cor. 16:17).

In 2 Cor. 8:18 a “brother” is mentioned who not only has praise for his zeal in the Gospel in all the congregations, but also “has been appointed by the congregations to travel with us with this gift,”

and in 8:23, Titus and our brethren are the congregations’ messengers and the glory of Christ.

In Phil. 4:15 Epaphroditus is called Paul’s brother and fellow laborer and fellow soldier, but the congregation’s (the Philippians’) messenger and minister to relieve Paul’s need.

Epaphras came to Paul, sent by the congregation in Colossae (Col. 4:12), perhaps also from Laodicea and Hierapolis.

From Thessalonica the Word sounded forth over Macedonia and Achaia and other places (1 Thess. 1:7).

In the Epistle to Titus (3:12) Paul says that he will send him Artemas or Tychicus, and asks him to help Zenas (the lawyer) and Apollos on their way, “that nothing be lacking unto them.”

Paul also sends the slave Onesimus to Philemon, and

Peter sends a letter by Silvanus (1 Pet. 5:12, 13).

2. Letters

were another excellent bond between the congregations, and testimony from the Word of God shows us that they also were diligently employed.

There were first letters from congregation to congregation. In Acts 15:22 we have in full a letter cited from the congregation in Jerusalem to the Gentile-Christian congregations which had already by that time been founded round about.

In several places letters of commendation from congregations and to congregations are mentioned, for example 2 Cor. 3:1 and 2 Thess. 2:2; in Acts 18:27 it is expressly communicated that the congregation (the brethren) in Ephesus sent a letter of commendation with Apollos to the congregations (the disciples) in Achaia and Corinth; and in 1 Cor. 16:3 Paul says that when he comes to Corinth, he will send their delegates with letters to Jerusalem in order to deliver the collected contributions; and so in other places.

Secondly, we have the significant collection of New Testament letters, besides hints of other letters not preserved, which more than anything else shows what an extraordinarily lively connection must have taken place between the congregations in this apostolic age. Paul’s greetings in Romans 16 furnish an excellent illustration of this relation, as also does the sixteenth chapter of First Corinthians. It almost seems as though letter-writing replaced what in later times has been carried out by Christian literature.

That the congregations also, where occasion was given, held

3. Joint Meetings

is beyond doubt enough, although such are expressly mentioned only in a couple of places, namely Acts 15:1 ff. and 20:16 ff. More frequently, however, there is mention of

4. Mutual Aid

as one of the means by which the congregations were bound more closely to one another. Thus the congregation in Antioch sent aid to those in need in Jerusalem (Acts 11:29).

Paul had undertaken, as a part of his missionary activity among the Gentiles, to remember the poor in Jerusalem with aid (Gal. 2:10). In accordance with this, he also arranged a general collection for those in need in Judea and obtained delegates, chosen by the congregations, to accompany him in order to deliver the aid (2 Cor., chs. 8 and 9; 1 Cor. 16:1).

What intimate connection was in this way brought about between the congregations, Paul describes in 2 Cor. 9:12 ff.: “For the service rendered by this support not only supplies the needs of the saints, but also bears rich fruit through many thanksgivings to God, while they, because of the proved spirit which this support shows, glorify God for the obedience of your confession to the Gospel of Christ and for your sincere liberality toward them and toward all, while they also in their prayer for you long after you because of God’s exceeding grace toward you. But thanks be to God for His unspeakable gift!”

This is the true Christian spirit of fellowship.

Thus it also manifested itself in the aid which the congregation in Philippi sent Paul several times, while he labored among congregations from whom he would receive no wages (Phil. 4:15–16).


% Page 84



5. Mutual Prayer

was thus also one of the most powerful means by which the spiritual bond and the consciousness of unity in Christ were strengthened and maintained.

6. Shared Orderings

we also find in the New Testament several traces of, as testimonies to, a vigorous and living sense of unity among the free congregations.

In Acts 14:23 it is related how Paul and Barnabas, on the return journey, traveled through the congregations which they had founded on the first missionary journey and had them choose elders in every congregation. In 1 Cor. 11:2 Paul praises the Corinthians because they hold fast the ordinances which he delivered to them, and after having given some most necessary admonitions concerning the Supper, he closes the chapter with the words: "The rest I will set in order when I come."

In the First Epistle to Timothy he gives determinations with regard to the choosing of elders or bishops and deacons, "in order," as he says, "that you may know how one ought to conduct oneself in God's House, which is the congregation of the living God." In the fifth chapter of the same epistle he gives rules for the choice of "widows," as likewise concerning the relation to the elders. Other examples may be adduced; but this is enough to show that fellowship in shared order is not of slight importance for the connection between free congregations.

\subsection{Free Fellowship between Free Congregations}

We have thus shown that there was indeed an exceedingly lively and extensive free fellowship between the free New Testament congregations, and that, despite differences of place and nationality, by their mutual connection they gave a powerful expression to the holy Church or congregation (Eph. 2:12ff. and 4:15-16).

The question now is this:
Are we the same congregation as then? Are we of the same kind and essence as they? For if that is so, then such a fellowship must by God's grace also be possible among us.

I hold that there is only one congregation, the same both in Heaven and on earth, and that if therefore there is at all a congregation of God among us, then it is the same as in the time of the apostles.

It was not the apostles' distinguished fear of God that brought forth the congregation. They had to wait until God's Spirit was poured out upon them; only then did there come to be a congregation.

It is God's Spirit who brings forth the congregation, whether among them or among us. There may be great difference in outward circumstances, as in inward development and gifts, but the Spirit is the same.

There is difference among human beings; they may be strong or weak, great or small, sick or healthy, children or old, dying or newborn: yet they all remain human beings; indeed, a dead human being is still a human being and not an animal, and what it lacks is not human nature, but life.

So also with the congregation. It is God's Spirit who, through the Word and the Sacraments, brings forth the congregation in believing human hearts now as then. And the Word and the Sacraments we have, and God's Spirit still works through them to this very day and is received in faith by penitent hearts, so that the congregation, with its salt and its leaven, does in truth exist among us.

So far as that is concerned, there is nothing to hinder free fellowship.

It is true, there is also difference both in inward conditions and in outward circumstances, some to the harm and others to the benefit of a free fellowship.

In one respect there is great likeness between us and the first congregations, namely this, that we too here in the land must make, as it were, a new beginning.

But in another respect there is just as great a difference, namely this, that among us, in a so-called Christian country, especially since we have come out of a State Church, where by the outward civil arrangements themselves one is, so to speak, forced into the congregation, there is almost no difference between the congregation and the world. This has worked incalculable harm for the congregation also here in the land, without thereby abolishing it, although it may with justice, and indeed in far too many cases, be said to our congregations as John says to the congregation in Sardis: "I know your works, that you have a name for being alive, and yet you are dead."

On the other hand, Christianity has not through these many hundreds of years been without influence upon the world either. Through so-called civilization it has to some degree permeated the world, its ways of thinking and its legislation, with the fundamental principles of Christian morality. At the same time as outward persecution has therefore almost wholly ceased, Christian and worldly morality have unfortunately often been mingled together.

At the same time, the means of outward connection have also, especially in the last century, been developed and refined to a high degree. It is now easier for one congregation to send a messenger of the Gospel to the end of the world than it was for Paul and other messengers of the congregations to come from the one end of the Mediterranean to the other.

Because of the development of the art of printing, it is now possible to place into the hand of every single member of the congregation the book which even Luther, only four hundred years ago, regarded it almost as an inestimable treasure to be able personally to own. The Bible, the source of life and faith both for the congregations and for individual Christians, can and ought now to be found in every house and can thus form the inward living bond between free congregations in a greater extent than even in the apostolic age.

A no less significant bond between our Lutheran congregations is formed by the Catechism. After long controversies it has now no doubt been acknowledged by the greatest part of our Church here as the congregation's own possession, the greatest blessing which the Lutheran Reformation has bestowed upon the Norwegian Church, a pure and simple presentation of the order of salvation grounded upon the Word and upon the Word alone. With the Lutheran Catechism in its hand and in its heart, the congregation can independently decide what is pure doctrine according to God's Word in contrast to all sorts of hair-splitting clerical babblings and theological speculations. By virtue of the Catechism, the congregation can therefore defend itself against the spirit-devouring theological controversies which both immediately after the apostolic age and in the Lutheran Church have again and again undermined all living Christianity. In the Catechism, therefore, our congregations in this land possess not only an advantage even over against the apostolic age, but also a spiritual basis of unity which, rightly used in connection with the awakening of new life, will form a stronger basis of fellowship between free Lutheran congregations than the most priestly systems of social organization.

But in other respects we stand immeasurably far behind the congregations of the New Testament age. The three ecclesial principles of unity mentioned above have in the course of centuries set their almost indelible carnal mark upon all church bodies. Among us too, the congregations, and with them often also believing pastors, have become so bound in chains by the so-called social organizations, by church politics and clerical regimes, that the congregation has almost forgotten that it exists for any other reason than purely outward ecclesiastical ends. Naturally the congregations have become unable to stir themselves for the high purposes for which the congregation has been set in the world, the mutual coming to one another and help toward the spread of God's Kingdom, which Paul so movingly depicts in Eph. 4:11-16, but which in most church bodies has been made into a mere clergy matter. Who has heard of congregations as such meeting together for mutual edification and encouragement in a State Church or a clerical society? And yet behold what intimate brotherly fellowship took place between the congregations in Christianity's first century!

But worse still, and essentially as a consequence of the congregations' state of bondage under the yoke of the clergy, a great multitude of our congregations are dead and thus altogether unfit to carry out the holy task God has entrusted to them. It is with this before his eyes that the apostle cries out in the letter to the congregation in Sardis: "Be watchful, and strengthen what remains, which is about to die!"

There is therefore, as said at the beginning, one presupposition indispensably necessary if the congregations are to be able to establish a free and living fellowship of unity among themselves, namely this, that they themselves become made alive, set free in Christ.

Praise be to the Lord! It is precisely this presupposition that is beginning to be realized among us, and which on the whole makes it possible to speak with hope of "free fellowship between free congregations," or, what is the same thing, the building up of a true and living Free Church among us. Tribulations and sufferings have in this respect wrought wonders among us as in the first Christian age; from many of our congregations in all church bodies there is heard of awakenings; the Lord has raised up particular instruments for this work of help in distress and will raise up more. The song of the turtledove has once more begun to sound, and the congregations to awaken to spiritual self-consciousness; and however weak and frail many may be, like bruised reeds and smoldering wicks, and however anxious and troubled believing pastors round about may feel, there is nevertheless now felt both obligation and need for a free and living fellowship between the awakening congregations, and for a fruitful cooperation, unhampered by social bonds, for the living building up of God's Kingdom both among ourselves mutually and among the heathen.

Small and groping attempts we have already made, and the Lord has not withheld His blessing. What now matters is only in childlike faith to set the course straight for the goal, in all things to let oneself be guided by the star of Bethlehem, God's Word and the first congregation. To the right and to the left we see shipwreck and ruins, congregations and peoples fallen into spiritual bondage under the spirit-extinguishing scepter of clerical rule. In the New Testament age we still see the congregation filled with God's Spirit and unbound by the later episcopate and state rule; Christ's freedom and an intimate fellowship of the Spirit united these congregations into one spiritual body, without priestly rule and social chains.

Thither, then, we must steer. We have shown that our congregation, though in great frailty, is nevertheless the same congregation of Christ as in the time of the apostles; for there is only one. The free fellowship which was possible for the first congregations must therefore by God's grace and the mighty working of His Spirit also be possible among us. Only we must go to the work with faith in the power of God's Spirit alone, and not lose courage over the many difficulties, and because that seems new which is nevertheless as old as God's Word.

This lecture has already grown beyond the bounds for an essay intended for the annual report.
We can therefore only in all possible brevity indicate the means through which free congregations after the New Testament pattern might come into free, living, and intimate fellowship of the Spirit with one another. Most of these means are neither unknown nor untried.

For the sake of clarity we may divide them into two great branches:

1. Mutual Help.

It will be rendered in general by the congregations' coming into a closer, so to speak personal, relation to one another and visiting one another as one family visits and in many ways assists another.

First, by remembering one another in prayer, just as one person prays for another person. By emissaries or messengers, pastors or laymen, who, each with his particular gifts, are so to speak exchanged from congregation to congregation, so that one congregation might thereby make good what the other lacks. By letters to one another and joint gatherings for edification for those congregations that are near enough to one another for this.

Second, by meetings, conversation meetings, larger or smaller, for mutual instruction and admonition, and for a deeper penetration into God's Word and the Catechism. Joint meetings for children and special revival meetings between congregations according to that remarkable word of Paul in Rom. 1:11, where he says with respect to the congregation at Rome: "For I long to see you, that I may impart to you some spiritual gift, so that you may be strengthened, that is, that I together with you may be encouraged among you by the faith we share, both yours and mine."

2. Cooperation.

First and foremost through mission work, both at home and abroad.

Next, through works of mercy of every kind, deaconess homes, hospitals, children's homes, etc.

Third, through the spread of edifying and awakening literature, papers and books, with special regard to the congregation, its nature and work, so that the congregations might thereby have a constant spiritual meeting place, for the exchange of thoughts and mutual encouragement to watch and pray and endure in the conflict.

Finally, fourth, through schools for the education of pastors and teachers and for the preservation of an unadulterated Word of God, which Satan has often used the keenest scholars to distort. "For we do not wrestle against flesh or blood, but against principalities, against powers, against the rulers of the world who rule in the darkness of this age, against the spiritual host of wickedness under heaven." It is for these purposes no less than for pastoral work that a school for pastors, conducted in God's Spirit with every competence, is altogether indispensable, especially for a Free Church.

This is surely only a scattered indication of what possibilities and means are present and needed for a "free fellowship between free congregations."

Continued treatment and practical experience will by God's gracious assistance demonstrate that a Lutheran Free Church in the sense in which it has here been presented, and will no doubt be acknowledged by this meeting, is not any mere thought-experiment or an unattainable ideal, but truly the old, good church order given by God's Word itself, in contrast to all priestly rule, denominational tyranny, and the state-church system.


